% Standalone lecture script. Compile with: latexmk -pdf PPO.tex
\documentclass[11pt,a4paper]{article}
\usepackage[T1]{fontenc}
\usepackage[utf8]{inputenc}
\usepackage{newtxtext,newtxmath}
\usepackage[margin=23mm,top=22mm,bottom=23mm,headheight=15pt]{geometry}
\usepackage{microtype}
\usepackage{amsmath}
\usepackage{xcolor,booktabs,tabularx,enumitem}
\usepackage[most]{tcolorbox}
\usepackage{tikz,pgfplots}
\usepackage{fancyhdr,titlesec}
\usepackage{hyperref}
\pgfplotsset{compat=1.18}
\definecolor{ink}{HTML}{18374A}
\definecolor{accent}{HTML}{087F8C}
\definecolor{pale}{HTML}{EEF6F6}
\definecolor{muted}{HTML}{536572}
\definecolor{warm}{HTML}{BD6735}
\hypersetup{colorlinks=true,linkcolor=accent,urlcolor=accent,citecolor=accent,
  pdftitle={PPO},pdfauthor={RL II 2026}}
\pagestyle{fancy}
\fancyhf{}
\fancyhead[L]{\small\sffamily\color{muted}RL II 2026}
\fancyhead[R]{\small\sffamily\color{muted}PPO}
\fancyfoot[R]{\small\sffamily\color{muted}\thepage}
\renewcommand{\headrulewidth}{0.3pt}
\titleformat{\section}{\Large\sffamily\bfseries\color{ink}}{\thesection}{0.65em}{}
\titleformat{\subsection}{\normalsize\sffamily\bfseries\color{ink}}{}{0pt}{}
\titlespacing*{\section}{0pt}{0pt}{12pt}
\titlespacing*{\subsection}{0pt}{10pt}{4pt}
\setlength{\parindent}{0pt}
\setlength{\parskip}{6pt plus 1pt}
\setlist{leftmargin=1.5em,itemsep=4pt,topsep=4pt}
\newtcolorbox{idea}[1]{enhanced,breakable,colback=pale,colframe=accent,
  boxrule=0pt,borderline west={2pt}{0pt}{accent},arc=0pt,
  left=9pt,right=9pt,top=7pt,bottom=7pt,
  title={#1},coltitle=ink,fonttitle=\sffamily\bfseries,
  attach title to upper={\par\smallskip},before skip=9pt,after skip=9pt}
\newcommand{\E}{\mathbb{E}}
\newcommand{\old}{\mathrm{old}}
\newcommand{\clip}{\operatorname{clip}}
\newcommand{\KL}{\operatorname{KL}}
\newcommand{\sg}{\operatorname{sg}}
\newcommand{\Ahat}{\widehat A}
\newcommand{\Vhat}{\widehat V}
\newcommand{\Lclip}{L^{\mathrm{clip}}}
\newcommand{\Lsur}{L^{\mathrm{sur}}}
\newcommand{\nexttopic}{\clearpage}
\begin{document}

{\sffamily\color{accent}\small LECTURE SCRIPT\par}
{\sffamily\bfseries\color{ink}\fontsize{27}{31}\selectfont PPO\par}
\vspace{9pt}

\section{PG Recap and REINFORCE}
Consider an episodic Markov decision process with horizon $H$ ($H=\infty$ works 
analogously, as seen in the RL lecture) and a differentiable
stochastic policy $\pi_\theta(a\,; s)$. At time $t$, the agent observes $S_t$,
chooses $A_t$, receives $R_t$, and moves to $S_{t+1}$. For a finite-horizon
problem, the state includes time, so we can write $V^\pi(S_t)$ without a
separate time index. Episodes that finish early can be padded by an absorbing
state with zero reward. The objective and the reward-to-go are
\begin{equation}\label{eq:objective}
 J(\theta)=\E_{\pi_\theta}\!\left[\sum_{t=0}^{H}\gamma^t R_t\right],
 \qquad G_t=\sum_{k=t}^{H-1}\gamma^{k-t}R_k,
 \qquad 0<\gamma\leq 1.
\end{equation}
Policy Gradient algorithms hinge on the policy gradient theorem, which gives
the gradient of the objective in terms of the policy's score function and the
reward-to-go:
\begin{equation}\label{eq:reinforce}
 \nabla J(\theta)=\E_{\pi_\theta}\!\left[
 \sum_{t=0}^{H-1}\gamma^t\nabla_\theta\log\pi_\theta(A_t\,; S_t)G_t\right].
\end{equation}
The discount factors in the front and inside the reward-to-go have different 
roles: $G_t$ discounts relative to \emph{time $t$}, while the outer $\gamma^t$ 
discounts relative to the \emph{start of the episode}.

\subsection{The REINFORCE update}
Collect $B$ independent episodes under the current policy and form
\begin{equation}\label{eq:reinforce-update}
 \widehat \nabla_\theta J(\theta)=\frac1B\sum_{i=1}^B\sum_{t=0}^{H-1}
 \gamma^t\nabla_\theta\log\pi_\theta(a_t^i\,; s_t^i)G_t^i,
 \qquad \theta\leftarrow\theta+\alpha\widehat \nabla_\theta J(\theta).
\end{equation}
This is stochastic gradient \emph{ascent}: actions followed by high returns
receive a positive push in log probability. Under the usual regularity
assumptions, the estimator is unbiased for~\eqref{eq:objective}.

\subsection{Algorithm 1: Vanilla REINFORCE}
{\small
\begin{tabbing}
\hspace{1.3em}\=\hspace{1.3em}\=\kill
\textbf{Initialize} policy parameters $\theta$; choose batch size $B$ and step size $\alpha$.\\
\textbf{repeat}\\
\> Collect $B$ complete episodes using the current policy $\pi_\theta$.\\
\> \textbf{for each} episode $i$ \textbf{do}\\
\>\> Set $G_H^i\leftarrow0$; compute $G_t^i\leftarrow r_t^i+\gamma G_{t+1}^i$ backwards.\\
\> Form the batch gradient $\widehat \nabla_\theta J(\theta)$ in~\eqref{eq:reinforce-update} using these returns.\\
\> Update $\theta\leftarrow\theta+\alpha\widehat \nabla_\theta J(\theta)$; discard the batch.\\
\textbf{until} the training budget is exhausted.
\end{tabbing}}

\begin{idea}{Two problems determine the route to PPO}
\textbf{Variance:} $G_t$ contains randomness from many later decisions and
transitions. Long episodes make gradient estimates very noisy.\par
\textbf{Sample efficiency:} after an update, the rollouts are discarded. In real-life examples,
what is expensive (real-time execution, damages from bad rollouts,...) is the cost of generating new data.
\end{idea}

\nexttopic
\section{A2C: Advantage Actor-Critic}
Define the state-action value function and the state value function of a policy by
\begin{equation}
 Q^\pi_t(s,a)=\E_\pi[G_t\mid S_t=s,A_t=a],\qquad
 V^\pi_t(s)=\E_{a\sim\pi(\cdot\mid s)}[Q^\pi_t(s,a)].
\end{equation}
Starting from the policy gradient theorem, we obtain
\begin{align*}
 \nabla J(\theta)
 &=\sum_{t=0}^{H-1}\gamma^t\E_{\pi_\theta}\!\left[
 \nabla_\theta\log\pi_\theta(A_t\mid S_t)G_t\right]
 =\sum_{t=0}^{H-1}\gamma^t\E_{\pi_\theta}\!\left[
 \nabla_\theta\log\pi_\theta(A_t\mid S_t)Q^{\pi_\theta}_t(S_t,A_t)\right]\\
 &=\sum_{t=0}^{H-1}\gamma^t\E_{\pi_\theta}\!\left[
 \nabla_\theta\log\pi_\theta(A_t\mid S_t)
 \bigl(Q^{\pi_\theta}_t(S_t,A_t)-b(S_t)\bigr)\right]\\
 &=
 \sum_{t=0}^{H-1}\gamma^t\E_{\pi_\theta}\!\Bigl[
 \nabla_\theta\log\pi_\theta(A_t\mid S_t)
 \underbrace{\bigl(Q^{\pi_\theta}_t(S_t,A_t)-V^{\pi_\theta}_t(S_t)\bigr)}_
 {\displaystyle :=A^{\pi_\theta}_t(S_t,A_t)}\Bigr].
\end{align*}
Here, we used the tower property in the first equation, the baseline trick from the
RL lecture in the second, and the value function as baseline in the third.
Conditioning removes future-return sampling noise from each gradient term.
Subtracting the value baseline then measures an action relative to the policy's
usual return in the same state, removing variation due to how favorable that
state is. This typically reduces gradient variance without changing its mean;
the value baseline is not necessarily the variance-minimizing baseline. The
differences $A^{\pi_\theta}_t(S_t,A_t)$ are called \emph{advantages}.

\subsection{Actor-critic Recap and A2C}
Actor-critic methods alternate two tasks: \emph{policy evaluation} estimates
how good the current policy is (the critic), and \emph{policy improvement}
changes the policy using that evaluation (the actor). In the RL lecture, we had 
model-based policy iteration as an example. We alternate between evaluating the 
current policy by solving its Bellman equations (critic; policy evaluation), then 
make the policy greedy with respect to that value (actor; policy improvement):
\begin{equation}
 \begin{gathered}
 V^\pi(s)=\sum_a\pi(a\mid s)\sum_{s'}P(s'\mid s,a)
 \bigl[r(s,a,s')+\gamma V^\pi(s')\bigr],\\[3pt]
 \pi_{\mathrm{new}}(s)\in\arg\max_a\sum_{s'}P(s'\mid s,a)
 \bigl[r(s,a,s')+\gamma V^\pi(s')\bigr].
 \end{gathered}
\end{equation}
In sample-based actor-critic methods, a learned critic replaces exact evaluation
and a general policy gradient step (e.g. through policy gradients) replaces the greedy update.
A2C learns $V^\pi$ by TD learning to evaluate the current policy: this is the
\emph{critic}. The policy $\pi_\theta$ takes stochastic gradient ascent steps
weighted by estimated advantages: this is the \emph{actor}. Now we only need a method to 
estimate the advantages from the critic's $V^\pi$.

\subsection{Generalized advantage estimation (GAE)}
\begingroup
\setlength{\abovedisplayskip}{5pt}
\setlength{\belowdisplayskip}{5pt}
\setlength{\abovedisplayshortskip}{3pt}
\setlength{\belowdisplayshortskip}{3pt}
For fixed critic predictions $\Vhat_t$, the one-step and $k$-step TD errors are
\begin{equation}\label{eq:td}
 \delta_t=R_t+\gamma\Vhat_{t+1}-\Vhat_t,
 \qquad
 \Ahat_t^{(k)}=\sum_{\ell=0}^{k-1}\gamma^\ell R_{t+\ell}
 +\gamma^k\Vhat_{t+k}-\Vhat_t\overset{\text{Telescope}}{=}\sum_{l=0}^{k-1}\gamma^l\delta_{t+l}.
\end{equation}
Short horizons typically reduce variance but rely on a potentially biased
bootstrap; longer horizons use more noisy rewards. GAE balances these effects
with $\lambda\in[0,1]$. For $m$ remaining transitions, define
\begin{equation}\label{eq:gae}
 \Ahat_t^{\mathrm{GAE}}=(1-\lambda)\sum_{k=1}^{m-1}\lambda^{k-1}\Ahat_t^{(k)}
 +\lambda^{m-1}\Ahat_t^{(m)}
 =\sum_{l=0}^{m-1}(\gamma\lambda)^l\delta_{t+l}.
\end{equation}
\endgroup

\nexttopic
\section{Sample reuse}
To tackle the problem of sample efficiency, we can reuse previously generated samples. 
Let $\pi_\old=\pi_{\theta_\old}$ generate a batch of trajectories. To take several policy updates
with these data, importance sampling (IS) changes the sampling distribution
from $\pi_\theta$ to $\pi_\old$. Assuming sufficient support, the gradient is
\begin{equation}\label{eq:exactis}
 \nabla J(\theta)=\sum_{t=0}^{H-1}\gamma^t\E_{\pi_\old}\!\left[
 \left(\prod_{j=0}^{t}\frac{\pi_\theta(A_j\mid S_j)}{\pi_\old(A_j\mid S_j)}\right)
 \nabla_\theta\log\pi_\theta(A_t\mid S_t)A_t^{\pi_\theta}(S_t,A_t)\right].
\end{equation}
The environment probabilities cancel. Earlier action factors correct the
state distribution; the last factor corrects the action distribution at $S_t$.
Here IS permits sample reuse, rather than itself reducing variance.

\begin{idea}{Three problems with exact sample reuse}
\textbf{Products and fractions:} Fractions of probabilities, especially if 
combined with long products can become very large, giving a
high-variance gradient estimate.\par
\textbf{Advantage estimation:} the formula
requires advantages for the current $\pi_\theta$, which can only be estimated with
samples from the current policy.\par
\textbf{Correlated rollouts:} neighboring transitions share history; using
whole rollouts as optimization units gives strongly correlated contributions.
\end{idea}

\subsection{How to solve the problems: local weighting, old advantages, random transitions}
Keep only the last factor of the product and introduce a trust region
$C_t(\theta_\old)$: the action probability has changed too far in the favorable
direction while we still work with stale data. The gate $c_t(\theta)=\mathbf{1}_{\{\theta\notin C_t(\theta_\old)\}}$
suppresses those contributions, limiting how long we trust the stale data. Replace the current advantage by the old one, then sample
individual transitions uniformly. Write
$r_t(\theta)=\tfrac{\pi_\theta(A_t\,; S_t)}{\pi_\old(A_t\,; S_t)}$ and abbreviate
$A_t^\old=A_t^{\pi_\old}(S_t,A_t)$. For $U\sim\mathrm{Unif}\{0,\ldots,H-1\}$,
\begin{align*}
 \nabla J(\theta)
 &\approx\sum_{t=0}^{H-1}\gamma^t\E_{\pi_\old}\!\left[
 \frac{\pi_\theta(A_t\mid S_t)}{\pi_\old(A_t\mid S_t)}c_t(\theta)
 \nabla_\theta\log\pi_\theta(A_t\mid S_t)A_t^\old\right]\\
 &=\sum_{t=0}^{H-1}\gamma^t\E_{\pi_\old}\!\left[
 \frac{\nabla_\theta\pi_\theta(A_t\mid S_t)}{\pi_\old(A_t\mid S_t)}
 c_t(\theta)A_t^\old\right]\\
 &=H\,\E_{U,\pi_\old}\!\left[
 \gamma^U\nabla_\theta r_U(\theta)c_U(\theta)A_U^\old\right].
\end{align*}
The second line uses $\pi\nabla\log\pi=\nabla\pi$; the third rewrites a sum as a
uniform expectation. In a buffer, replace $A_t^\old$ by its fixed GAE estimate
$\Ahat_t$ and mix transitions from all rollouts in minibatches. Shuffling mixes
correlated observations; it does not make the original data independent.

We deliberately trade \textbf{bias in the gradient} for lower variance and the
ability to reuse expensive samples. Keeping the policy near the version that
collected the data makes old advantages more relevant. The precise gate is the
PPO clipping rule below. A more standard motivation comes from trust-region
policy optimization, which we will likely encounter in a later lecture.

Common PPO losses omit the outer $\gamma^t$. To retain
it, either weight uniform buffer samples by $\gamma^t$, or draw transitions with
probability proportional to $\gamma^t$ (up to a constant normalization), as in
the second lecture's homework. Below, $w_t=\gamma^t$ denotes the first option;
$w_t=1$ denotes the usual uniform PPO convention. Here $t$ is episode time.

\nexttopic
\section{PPO: Proximal Policy Optimization}
\subsection{The choice of trust region}
Initially $r_t(\theta)=r_t(\theta_\old)=1$. The sign of the stored advantage determines
which change is favorable.
\begin{itemize}
 \item If $\Ahat_t>0$, increase the probability of $A_t$. Once
 $r_t>1+\epsilon$, stop rewarding further increases. If the probability instead
 decreases, keep the gradient that pushes it back up.
 \item If $\Ahat_t<0$, decrease the probability of $A_t$. Once
 $r_t<1-\epsilon$, stop rewarding further decreases. If the probability instead
 increases, keep the gradient that pushes it back down.
\end{itemize}
The asymmetric clipping is so that information from opposite advantages is not lost.

\subsection{The PPO loss}
The derivation gives an improvement direction in the form of a gradient. The neural 
networks used for deep reinforcement learning cannot generally
realize arbitrary changes to individual action probabilities, or be optimized
in one exact step. Instead, we specify a differentiable loss and let a numerical
optimizer take parameter steps on minibatches. Its computed gradients are
those of the sampled loss, rather than the exact expected return. We will
study network training and these optimizers (especially for the critic using TD-Learning) more closely in later lectures.
We obtain the classical PPO loss:

\begin{idea}{The PPO Loss}
\begin{equation}\label{eq:clip}
 \Lclip(\theta)=\widehat\E_{U,\pi_\old}\!\left[w_U
 \min\!\left\{r_U(\theta)\Ahat_U,\;
 \clip\bigl(r_U(\theta),1-\epsilon,1+\epsilon\bigr)\Ahat_U\right\}\right].
\end{equation}
\end{idea}

\subsection{Algorithm 2: PPO}
{\small
\begin{tabbing}
\hspace{1.3em}\=\hspace{1.3em}\=\kill
\textbf{Initialize} policy parameters $\theta$ (actor), value function $V$ (critic);\\
\> \textbf{choose} batch size $B$, number of epochs $E$, and step size $\alpha$.\\
\textbf{repeat}\\
\> Collect $B$ complete episodes using the fixed policy $\pi_\theta$ and compute advantages backwards.\\
\> Draw $E$ minibatches of size $M$ from the buffer.\\
\> Alternately update actor $\theta\leftarrow\theta+\alpha\widehat\nabla_\theta J(\theta)$ and critic $V$ (TD learning) on each minibatch.\\
\textbf{until} the training budget is exhausted.
\end{tabbing}}

\subsection{A few differences in practice}
\begin{itemize}[itemsep=3pt]
 \item \textbf{Shared networks:} actor and critic may share parameters. Then
 use a combined actor and value-prediction loss and a joint optimizer step, rather than
 two separate updates as in the tabular standard algorithm.
 \item \textbf{Finite and infinite horizons:} a rollout is necessarily a finite segment. Bootstrap at nonterminal cutoffs from the final
 observation before any reset; use zero value at true termination.
 \item \textbf{Adam versus SGD:} implementations often use Adam (or other more sophisticated optimizers) in place of
 SGD steps; it adapts steps using gradient history.
 \item \textbf{Unequal rollout lengths:} episodes or collected segments may
 have different lengths. The buffer then contains $\sum_iT_i$ transitions;
 sample transitions, and compute GAE separately within each segment.
\item \textbf{Gamma} $\gamma$ is often treated as a hyperparameter in practice, allowing for tuning based on the specific task and environment. It is also common, as noted, to not use $\gamma$ in the outer weighting of the PPO loss.
\end{itemize}

\end{document}
